\chapter{Exercises and Worksheets}

This appendix provides practical exercises for developing DCF skills. Use these individually or in workshop settings.

\section{Exercise 1: Socratic Self-Assessment}

\textbf{Objective}: Evaluate your current interaction patterns with AI.

\textbf{Instructions}:
\begin{enumerate}
    \item Review your last 5 AI conversations (or start fresh with a technical question)
    \item For each exchange, classify your prompts:
    \begin{itemize}
        \item \textbf{Assertion}: ``Write me a function that...''
        \item \textbf{Question}: ``What are the tradeoffs between...''
        \item \textbf{Challenge}: ``Why did you choose X over Y?''
        \item \textbf{Meta}: ``Let's step back---what's the real problem here?''
    \end{itemize}
    \item Calculate the ratio of each type
    \item Identify patterns and set goals for more questioning
\end{enumerate}

\textbf{Reflection Questions}:
\begin{itemize}
    \item What percentage of your prompts were assertions vs. questions?
    \item When you asked questions, were they genuine inquiries or rhetorical?
    \item How often did you challenge the AI's initial response?
\end{itemize}

\section{Exercise 2: The Why Chain}

\textbf{Objective}: Practice depth of inquiry through sequential questioning.

\textbf{Instructions}:
\begin{enumerate}
    \item Start with any AI response you've recently received
    \item Ask ``Why?'' or ``How do you know?'' five times in sequence
    \item Document what you learn at each level
    \item Note when the AI reaches the limits of its knowledge
\end{enumerate}

\textbf{Example Chain}:
\begin{enumerate}
    \item AI: ``Use a HashMap for this data structure.''
    \item Human: ``Why a HashMap specifically?''
    \item AI: ``O(1) average lookup time.''
    \item Human: ``Why is O(1) lookup important here?''
    \item AI: ``Your access pattern is key-based retrieval.''
    \item Human: ``How do you know my access pattern?''
    \item AI: ``Based on the code you showed...''
    \item Human: ``What assumptions are you making?''
    \item AI: (Now we reach the interesting territory)
\end{enumerate}

\section{Exercise 3: Mirror Calibration}

\textbf{Objective}: Understand how input quality affects output quality.

\textbf{Instructions}:
\begin{enumerate}
    \item Choose a technical problem you understand well
    \item Craft three prompts of varying quality:
    \begin{itemize}
        \item \textbf{Low}: Vague, no context (``fix my code'')
        \item \textbf{Medium}: Some context, clear request
        \item \textbf{High}: Rich context, specific question, constraints stated
    \end{itemize}
    \item Submit each to the AI and compare responses
    \item Analyze the correlation between input and output quality
\end{enumerate}

\textbf{Scoring Rubric}:
\begin{table}[H]
\centering
\begin{tabular}{lp{8cm}}
\toprule
\textbf{Dimension} & \textbf{Assessment Questions} \\
\midrule
Relevance & Did the response address your actual need? \\
Accuracy & Was the technical content correct? \\
Completeness & Were important aspects covered? \\
Actionability & Could you implement the suggestion? \\
\bottomrule
\end{tabular}
\end{table}

\section{Exercise 4: Checkpoint Design}

\textbf{Objective}: Practice identifying effective approval points in a workflow.

\textbf{Instructions}:
\begin{enumerate}
    \item Choose a multi-step task (e.g., refactor a module, write documentation)
    \item Map out all the steps an agentic AI would take
    \item Identify which steps are:
    \begin{itemize}
        \item \textbf{Reversible}: Can be easily undone
        \item \textbf{Consequential}: Have significant impact
        \item \textbf{Uncertain}: Require human judgment
    \end{itemize}
    \item Design checkpoint placement that balances oversight with efficiency
    \item For each checkpoint, write 2-3 Socratic questions you would ask
\end{enumerate}

\section{Exercise 5: Recursive Refinement Practice}

\textbf{Objective}: Experience the iterative improvement process.

\textbf{Instructions}:
\begin{enumerate}
    \item Ask the AI to explain a concept you're learning
    \item Rate the initial explanation (1-10) and identify gaps
    \item Craft a follow-up question addressing the weakest point
    \item Repeat steps 2-3 until you rate the explanation 8+
    \item Count how many iterations were required
\end{enumerate}

\textbf{Tracking Sheet}:
\begin{table}[H]
\centering
\begin{tabular}{cccp{5cm}}
\toprule
\textbf{Iteration} & \textbf{Rating} & \textbf{Gap Identified} & \textbf{Follow-up Question} \\
\midrule
1 & & & \\
2 & & & \\
3 & & & \\
4 & & & \\
5 & & & \\
\bottomrule
\end{tabular}
\end{table}

\section{Exercise 6: Metacognitive Journaling}

\textbf{Objective}: Develop awareness of your own thinking during AI collaboration.

\textbf{Instructions}:
\begin{enumerate}
    \item During your next significant AI session, pause every 10 minutes
    \item Record:
    \begin{itemize}
        \item What was I trying to accomplish?
        \item What did I actually ask for?
        \item How did I feel about the AI's response?
        \item What am I assuming?
        \item What would I do differently?
    \end{itemize}
    \item After the session, review your notes for patterns
\end{enumerate}

\section{Exercise 7: Framework Comparison}

\textbf{Objective}: Understand DCF's position in the landscape of approaches.

\textbf{Instructions}:
\begin{enumerate}
    \item Choose a real task you need to complete
    \item Approach it using three different methodologies:
    \begin{itemize}
        \item Pure prompting (single detailed prompt)
        \item Agentic (let AI plan and execute autonomously)
        \item DCF (Socratic dialogue at checkpoints)
    \end{itemize}
    \item Compare results on: quality, time, learning, satisfaction
    \item Identify which approach suits which task types
\end{enumerate}

\section{Exercise 8: The Silence Test}\index{Silence Test}

\textbf{Objective}: Assess what you've internalized vs. what you still depend on AI for.

This exercise tests whether AI collaboration is building your capability or creating dependency---the core question of scaffolding theory.

\textbf{Instructions}:
\begin{enumerate}
    \item Choose a task you normally do with AI assistance
    \item Attempt it \emph{without} AI. Set a timer for 20 minutes.
    \item During the attempt, track:
    \begin{itemize}
        \item Where you got stuck
        \item What you wanted to ask AI for help with
        \item What you successfully did independently
    \end{itemize}
    \item After 20 minutes, use AI to complete the task
    \item Compare: What was genuinely difficult vs. what was habitual dependency?
\end{enumerate}

\textbf{Tracking Sheet}:
\begin{table}[H]
\centering
\begin{tabular}{p{4cm}p{4cm}p{4cm}}
\toprule
\textbf{Got Stuck On} & \textbf{Wanted AI For} & \textbf{Did Independently} \\
\midrule
& & \\
& & \\
& & \\
\bottomrule
\end{tabular}
\end{table}

\textbf{Reflection Questions}:
\begin{itemize}
    \item What should I practice doing without AI?
    \item What is AI genuinely augmenting vs. replacing?
    \item Am I building capability or dependency?
    \item If the same patterns appear repeatedly, what skill am I avoiding developing?
\end{itemize}

\textbf{The DCF Connection}: This exercise operationalizes the scaffolding principle---AI should help you reach capabilities you couldn't achieve alone, then fade as you internalize those capabilities. If you can't do \emph{anything} without AI, the scaffold has become a crutch.

\section{Workshop Format: DCF Introduction (90 minutes)}

\textbf{For facilitators introducing DCF to teams:}

\begin{table}[H]
\centering
\begin{tabular}{lll}
\toprule
\textbf{Time} & \textbf{Activity} & \textbf{Materials} \\
\midrule
0-10 min & Introduction to DCF concepts & Slides from Chapter 3 \\
10-25 min & Exercise 1: Self-Assessment & Individual work \\
25-35 min & Pair share: Assessment results & Discussion pairs \\
35-55 min & Exercise 2: Why Chain (live demo) & Shared screen \\
55-70 min & Exercise 3: Mirror Calibration & Small groups \\
70-85 min & Debrief and Q\&A & Full group \\
85-90 min & Commitment: One DCF practice to try & Individual \\
\bottomrule
\end{tabular}
\end{table}

